% !TEX root = ../Main.tex
\chapter{通用比赛规则}

\section{出发与触壁}

\state{出发方式}{
    \begin{itemize}
        \item \textbf{自由泳、蛙泳、蝶泳、个人混合泳}：必须从\textbf{出发台起跳}出发；
        \item \textbf{仰泳}：\textbf{在水中出发}（面对出发端、两手抓握手器）。
    \end{itemize}
}

\state{到达终点的触壁规则}{
    \begin{itemize}
        \item \textbf{自由泳、仰泳}：到达终点时可以\textbf{只用一只手触壁}；
        \item \textbf{蛙泳、蝶泳}：\textbf{必须使用双手同时触壁}。
    \end{itemize}
}

\section{转身规则}

\state{折返与转身}{
    \begin{itemize}
        \item 奥运会所有\textbf{距离在 50 米以上}的游泳比赛都\textbf{必须途中折返}；
        \item \textbf{自由泳、仰泳}转身：允许运动员使用\textbf{身体的任何部分}触及池壁，因此允许在水下转身后\textbf{用脚蹬池壁}；
        \item \textbf{蛙泳、蝶泳}转身：必须\textbf{双手同时触壁}（见各泳姿章节）；
        \item \textbf{个人混合泳特例}：由\textbf{仰泳转换泳姿到蛙泳}时，运动员\textbf{必须保持仰卧姿势触及池壁}。
    \end{itemize}
}

\section{电子计时与接力出发}

\state{电子计时系统}{
    \begin{itemize}
        \item 运动员的\textbf{比赛时间}和\textbf{触壁地点}都由\textbf{电子系统自动决定}；
        \item 出发时，\textbf{出发台上的压力板}记录数据；
        \item 每条泳道两边的池壁上装有\textbf{触摸板}，运动员触壁时被记录；
        \item 触摸板与出发台\textbf{互相连接}，可判定接力出发是否犯规。
    \end{itemize}
}

\state{接力出发犯规——0.03 秒规则}{
    \begin{itemize}
        \item 参加接力的运动员\textbf{在其队友触壁之前 0.03 秒}离开出发台即算\textbf{犯规}；
        \item 一旦犯规，整个队伍将被\textbf{自动取消比赛资格}。
    \end{itemize}
}

\rmkb{
    \textbf{易考点}：
    \begin{itemize}
        \item \textbf{只有仰泳}在\textbf{水中出发}，其他泳姿（含个人混合泳）均从出发台出发；
        \item \textbf{蛙泳 / 蝶泳}\textbf{双手同时触壁}；自由泳 / 仰泳\textbf{一只手即可}；
        \item 自由泳、仰泳转身允许\textbf{身体任何部分}触壁；
        \item 混合泳特例：\textbf{仰 $\to$ 蛙}转身要\textbf{保持仰卧触壁}；
        \item 接力\textbf{提前 0.03 秒}离台 $\to$ 取消资格。
    \end{itemize}
}
